\documentclass[12pt,a4paper]{article}

\usepackage{fontawesome5}
\providecommand{\faPencilAlt}{\faIcon{pencil-alt}}

\usepackage{styles/couleurs}
\usepackage{styles/mise_en_page}
\usepackage{styles/boites}
\usepackage{styles/maths}
\usepackage{styles/figures}

\usepackage{enumitem}
\usepackage{array}
\usepackage{tabularx}
\usepackage{amsmath}
\usepackage{amssymb}
\usepackage{tikz}
\usepackage{xparse}

\usetikzlibrary{arrows.meta,calc}

\setlist[itemize]{leftmargin=1.4em,itemsep=0.15em,topsep=0.2em}
\setlist[enumerate]{leftmargin=*,itemsep=0.25em,topsep=0.25em}

\newcommand{\lignepointillee}{%
  \noindent\dotfill\par
  \vspace{0.45em}
}

\newcommand{\lignesreponse}[1]{%
  \begin{tcolorbox}[
    enhanced,
    colback=white,
    colframe=NeutreC!40,
    boxrule=0.5pt,
    arc=1.5mm,
    left=3mm,
    right=3mm,
    top=2mm,
    bottom=1mm,
    before skip=2pt,
    after skip=2pt
  ]
    \foreach \i in {1,...,#1}{%
      \lignepointillee
    }
  \end{tcolorbox}
}

\newcommand{\titreseq}[1]{%
\begin{center}
  {\sffamily\large\bfseries\color{SavoirC}#1}
\end{center}
}

\begin{document}

\titreseq{Activité 4 -- Quotients de nombres relatifs}

\noindent
\begin{minipage}[t]{0.49\linewidth}

\section*{Partie A -- Facteur manquant}

Complète les égalités.

\[
(+6)\times\makebox[1.5cm]{\dotfill}=-42
\]


\[
(-5)\times\makebox[1.5cm]{\dotfill}=-35
\]

\[
(-8)\times\makebox[1.5cm]{\dotfill}=40
\]

\[
(+4)\times\makebox[1.5cm]{\dotfill}=18
\]

\end{minipage}\hfill
\begin{minipage}[t]{0.49\linewidth}

\section*{Partie B -- Relier à une division}

À partir de la partie A, complète.

\[
(-42)\div(+6)=\ldots
\]

\[
(-35)\div(-5)=\ldots
\]

\[
40\div(-8)=\ldots
\]

\[
18\div(+4)=\ldots
\]

\end{minipage}


\section*{Partie C -- Règle des signes du quotient}

Complète.

\begin{tcolorbox}[colback=SavoirC!5,colframe=SavoirC!60,boxrule=0.7pt,arc=2mm]
Le quotient de deux nombres de même signe est \dotfill

\vspace{0.4cm}

Le quotient de deux nombres de signes différents est \dotfill

\vspace{0.4cm}

Pour calculer un quotient, je détermine d'abord le \dotfill, puis je divise les distances à zéro.
\end{tcolorbox}

\section*{Partie D -- Calculatrice raisonnée}

Avant d'utiliser la calculatrice, écris le signe attendu et un ordre de grandeur.

\renewcommand{\arraystretch}{1.7}
\begin{tabularx}{\textwidth}{|>{\centering\arraybackslash}p{4cm}|>{\centering\arraybackslash}p{3cm}|>{\centering\arraybackslash}p{4cm}|>{\centering\arraybackslash}X|}
\hline
\textbf{Calcul} & \textbf{Signe attendu} & \textbf{Ordre de grandeur} & \textbf{Résultat}\\
\hline
\((-12{,}8)\div2\) & & & \\
\hline
\((-63)\div(-0{,}7)\) & & & \\
\hline
\(7{,}2\div(-5)\) & & & \\
\hline
\end{tabularx}

\newpage

\titreseq{Activité 5 -- Quotients de nombres relatifs}

\noindent
\begin{minipage}[t]{0.49\linewidth}

\section*{Partie A -- Facteur manquant}

Complète les égalités.

\[
(+6)\times\makebox[1.5cm]{\dotfill}=-42
\]


\[
(-5)\times\makebox[1.5cm]{\dotfill}=-35
\]

\[
(-8)\times\makebox[1.5cm]{\dotfill}=40
\]

\[
(+4)\times\makebox[1.5cm]{\dotfill}=18
\]

\end{minipage}\hfill
\begin{minipage}[t]{0.49\linewidth}

\section*{Partie B -- Relier à une division}

À partir de la partie A, complète.

\[
(-42)\div(+6)=\ldots
\]

\[
(-35)\div(-5)=\ldots
\]

\[
40\div(-8)=\ldots
\]

\[
18\div(+4)=\ldots
\]

\end{minipage}


\section*{Partie C -- Règle des signes du quotient}

Complète.

\begin{tcolorbox}[colback=SavoirC!5,colframe=SavoirC!60,boxrule=0.7pt,arc=2mm]
Le quotient de deux nombres de même signe est \dotfill

\vspace{0.4cm}

Le quotient de deux nombres de signes différents est \dotfill

\vspace{0.4cm}

Pour calculer un quotient, je détermine d'abord le \dotfill, puis je divise les distances à zéro.
\end{tcolorbox}

\section*{Partie D -- Calculatrice raisonnée}

Avant d'utiliser la calculatrice, écris le signe attendu et un ordre de grandeur.

\renewcommand{\arraystretch}{1.7}
\begin{tabularx}{\textwidth}{|>{\centering\arraybackslash}p{4cm}|>{\centering\arraybackslash}p{3cm}|>{\centering\arraybackslash}p{4cm}|>{\centering\arraybackslash}X|}
\hline
\textbf{Calcul} & \textbf{Signe attendu} & \textbf{Ordre de grandeur} & \textbf{Résultat}\\
\hline
\((-12{,}8)\div2\) & & & \\
\hline
\((-63)\div(-0{,}7)\) & & & \\
\hline
\(7{,}2\div(-5)\) & & & \\
\hline
\end{tabularx}



\end{document}
